\section{参数敏感性与模型评价}

\subsection{备用、风险权重与终端库存}

主策略采用备用系数$\kappa=1$、风险基数$\lambda_0=0.1$。二者作为事先固定的策略偏好，未利用2月至12月的实际费用寻优。表\ref{tab:sensitivity}完整保留增强参数带来费用上升的情况，以判断机制的实际边界。

备用系数从0.5提高至1.5时，问题三紧急量由14208.94 kWh降至11688.95 kWh，总费用由14023535.12元升至14067387.71元。更多参考备用增加了缓冲空间，但也限制了廉价电量的灵活使用，并可能导致更多承诺购电剩余。若运行目标仅是最小化本样本实际电费，较大的备用系数不一定合适。

风险基数取0时，总费用13927283.23元、紧急量14609.38 kWh；主策略费用14042726.52元、紧急量12685.74 kWh。主策略以约11.54万元额外支出减少约1923.64 kWh紧急需求，而且其日总费用CVaR95也高于零风险权重情形。原因在于优化惩罚的是预测场景中的紧急费用尾部，实际总费用还包含常规采购和调整费，且经验场景存在有限样本误差。不能将对一个随机损失分量的风险约束，直接解释为对全部实现成本尾部的单调改进。若不对紧急需求额外赋予偏好，应优先考虑零风险权重的经济型方案；本文保留既定主参数以展示明确的风险偏好，而不在评价期事后挑选最佳参数。

\begin{table}[htbp]\centering\small\caption{问题3风险与终端参数敏感性}
\label{tab:sensitivity}
\begin{tabularx}{\linewidth}{LRRRR}\toprule
策略 & 总费用/万元 & 紧急量/kWh & 日费CVaR95/元 & 年末电量/kWh\\\midrule
主策略：$\kappa=1,\lambda_0=0.1$ & 1404.27 & 12685.74 & 66381.40 & 6247.51\\
备用系数$\kappa=0.5$ & 1402.35 & 14208.94 & 66336.84 & 6247.51\\
备用系数$\kappa=1.5$ & 1406.74 & 11688.95 & 66461.26 & 6247.51\\
风险基数$\lambda_0=0$ & 1392.73 & 14609.38 & 65683.92 & 6245.12\\
风险基数$\lambda_0=0.2$ & 1416.64 & 11312.96 & 66878.46 & 6285.62\\
取消终端偏离惩罚 & 1403.32 & 18350.10 & 66351.31 & 1637.69\\
下调退款解释 & 1369.08 & 13430.61 & 64275.62 & 6234.10\\
\bottomrule\end{tabularx}\end{table}

取消终端惩罚后，费用降至14033228.53元，但紧急量升至18350.10 kWh，年末储电量降至1637.69 kWh，明显低于主策略6247.51 kWh。由于结存能量不同，不能把9497.99元费用下降全部认定为运营效率提升。对终端机制的比较应同时报告期初、期末电量与实际紧急需求；软目标只能抑制有限时域内过度放空的倾向，不能保证全天实际末态恰等于参考值。

\subsection{下调购电是否退款的解释敏感性}

主结果严格采用式\eqref{eq:settlement}的“计划全额付费加调整相关费用”口径。另一种可能解释是，下调时取消原计划对应电费，再支付50\%违约费，于是其净影响为$-0.5p_t[g_t-a_t]_+$。为避免把语义差异伪装成算法优势，另以这一退款解释重算问题三，结果仅用于敏感性比较，不替换附件5主结果。

退款解释下，总费用为13690758.70元，最终常规购电量少于原始计划量，累计下调1223373.21 kWh，调整净费用为$-160832.93$元。与主口径的费用差异约35.20万元，明显大于主策略在若干消融对照中的改善。因此结算语义对结果具有实质影响。该敏感性说明所报数值依赖明示的结算假设，不能在不同退款规则之间直接比较策略优劣。

\subsection{模型特点与适用边界}

本文的主要特点是维持一致的物理、信息与结算结构。各问使用同一储电状态方程和执行器，区别仅在预测来源、场景维度与计划更新时间。整日联合残差保留相关误差，避免逐时独立抽样低估持续预测偏差；参考备用把预测误差转换为储能量级，便于解释其为何抑制过早放电。所有费用以实际执行和实际电价结算，使优化参考成本与真实支付结果能够分开比较。

物理可行性通过逐时能量守恒、储电边界、充放电上限、互斥与跨日连续性计算检查。执行器每次先使用物理可行储能，再紧急补足负荷，所以即使场景没有覆盖实际误差，也不会把未满足负荷隐藏为零购电。外网容量未受限制是这一性质的前提，若实际存在接口限额，需进一步引入网络约束和失供电代价。

仍有四方面限制。第一，模型仅使用一个年度，且最初只有一个月用于选择预测器，难以据此声称跨年份泛化。第二，两阶段模型固定各场景共享的储能参考轨迹，实际再由反馈规则修正，没有求解任意历史依赖控制策略。第三，滚动时域只延伸到当日末端，跨日峰谷变化仅通过软终端目标近似。第四，容量、效率和响应速度按题目常数处理，未引入寿命退化、爬坡、交易次数或真实电网容量，因此结果适用于题目给定边界，不能无修改用于实际设备运行。

\section{结论}

确定性条件下，线性储能联合调度以35126.95元完成单日供电，节费26.90\%；当预测误差存在时，按历史残差配置购电和储能参考轨迹，比单纯围绕均值套利显著减少五倍补购造成的费用。规定时点的滚动更新在固定电价下具有可观的应急需求削减效果；电价同时波动时，总费用改善较小，主要表现为较低的紧急电量和较低的高费用日期成本。

风险机制应作为可选择的偏好参数，不能预设其总能降低全部成本。以实际电费为唯一目标与希望更少依赖紧急购电，是两种应明确区分的运行偏好。完整的预测、求解、执行与结算代码以及全部334天策略随支撑材料保存，便于根据最终采用的偏好和结算解释复算。
